\documentclass[11pt]{article}

\usepackage[T1]{fontenc}
\usepackage[utf8]{inputenc}
\IfFileExists{libertinus.sty}{\usepackage{libertinus}}{\usepackage{lmodern}}
\usepackage{microtype}
\usepackage[a4paper,margin=28mm,headheight=14pt]{geometry}
\usepackage{amsmath,amssymb,mathtools,amsthm}
\usepackage{bm}
\usepackage{booktabs,longtable,tabularx,array}
\usepackage{enumitem}
\usepackage{xcolor}
\usepackage{fancyhdr}
\usepackage{hyperref}
\usepackage[nameinlink,noabbrev]{cleveref}

\hypersetup{
  colorlinks=true,
  linkcolor=blue!45!black,
  citecolor=green!35!black,
  urlcolor=blue!55!black,
  pdftitle={q-Binomial Richardson Acceleration of Geometric Sinc Products},
  pdfauthor={Frontier research report for the ProveIt Fabius Function project}
}
\urlstyle{tt}
\setlength{\emergencystretch}{3em}

\pagestyle{fancy}
\fancyhf{}
\lhead{Geometric sinc products and Rvachev splines}
\rhead{Frontier report}
\cfoot{\thepage}

\newtheorem{theorem}{Theorem}[section]
\newtheorem{proposition}[theorem]{Proposition}
\newtheorem{corollary}[theorem]{Corollary}
\newtheorem{lemma}[theorem]{Lemma}
\newtheorem{conjecture}[theorem]{Conjecture}
\theoremstyle{definition}
\newtheorem{definition}[theorem]{Definition}
\newtheorem{example}[theorem]{Example}
\newtheorem{algorithm}[theorem]{Algorithm}
\theoremstyle{remark}
\newtheorem{remark}[theorem]{Remark}

\newcommand{\sinc}{\operatorname{sinc}}
\newcommand{\up}{\operatorname{up}}
\newcommand{\e}{\mathrm{e}}
\newcommand{\ii}{\mathrm{i}}
\newcommand{\R}{\mathbb{R}}
\newcommand{\C}{\mathbb{C}}
\newcommand{\N}{\mathbb{N}}
\newcommand{\D}{\mathrm{D}}
\newcommand{\norm}[1]{\left\lVert #1\right\rVert}
\newcommand{\abs}[1]{\left\lvert #1\right\rvert}
\newcommand{\qbinom}[2]{\genfrac{[}{]}{0pt}{}{#1}{#2}_{q}}
\newcommand{\repo}{\textsc{ProveIt}}
\DeclareMathOperator{\supp}{supp}
\DeclareMathOperator{\sgn}{sgn}

\title{\bfseries q-Binomial Richardson Acceleration of Geometric Sinc Products\\[2mm]
\large Bernoulli--Bell tail expansions, moment-matched Rvachev splines,\\
and interior inverse-Fabius asymptotics}
\author{Frontier research report for the \repo{} Fabius Function project}
\date{27 August 2026}

\begin{document}
\maketitle

\begin{abstract}
The Fourier transform of Rvachev's compactly supported up-function is the geometric
sinc product
\[
   P(\xi)=\prod_{\nu\geq 1}\sinc(\xi/2^\nu).
\]
The repository corpus develops its large-frequency decay, its connection with the
Fabius function and the signed Thue--Morse sequence, dyadic exact values,
small-argument and inverse asymptotics, and several Richardson-type limits.
This report studies a different asymptotic variable: the truncation level in the
finite product
\(P_n(\xi)=\prod_{\nu=1}^n\sinc(\xi/2^\nu)\).

For a general geometric base \(b>1\), with \(q=b^{-2}\), we prove, in the
disc \(\lvert q^n\xi^2\rvert<\pi^2b^2\), the total exact analytic factorization
\[
 P_{b,n}(\xi)
 =P_b(\xi)\exp\!\left(\sum_{k\geq1}c_k(b)\,\xi^{2k}q^{kn}\right)
 =P_b(\xi)\sum_{r\geq0}\alpha_r(b)\,\xi^{2r}q^{rn}.
\]
Off the zero set of \(P_b\), division gives the corresponding quotient form.
The cumulants \(c_k(b)\) are explicit Bernoulli-number combinations and the
coefficients \(\alpha_r(b)\) are complete Bell polynomials. Lagrange evaluation
at the geometric nodes \(1,q,\ldots,q^m\) then gives a universal extrapolant
\[
 \mathcal R_{m,n}P_b=\sum_{j=0}^m w_j^{(m)}P_{b,n+j},
 \qquad
 w_j^{(m)}=\frac{(-1)^{m-j}q^{(m-j)(m-j+1)/2}}
 {(q;q)_j(q;q)_{m-j}}.
\]
On the same disc and off the zero set of \(P_b\), its complete error series has
the closed q-binomial form
\[
 \frac{\mathcal R_{m,n}P_b(\xi)}{P_b(\xi)}-1
 =(-1)^m q^{\binom{m+1}{2}}
 \sum_{r\geq m+1}\qbinom{r-1}{m}\alpha_r(b)\xi^{2r}q^{rn}.
\]
Consequences include an exact sign law, non-asymptotic divided-difference bounds,
a rigorous a posteriori error certificate, and a uniformly bounded stability
constant \(\prod_{r=1}^m(1+q^r)/(1-q^r)\). For \(q=1/4\) this constant stays
below \(1.970\), even as the order tends to infinity.

After inverse Fourier transformation, the same filter combines consecutive finite
box splines. It matches the first \(2m\) moments of the limiting up-density
\emph{exactly at every truncation level}, and converges in every \(C^s\)-norm at
rate \(q^{(m+1)n}\). The first error profile is an explicit even derivative of
the up-function. Translating to the Fabius function gives the same accelerated
approximation, while inversion on compact interior intervals yields an explicit
first correction for the inverse Fabius function. These statements are proved
here and checked symbolically and numerically by the accompanying Python program.
They are described as \emph{repository-new}: the claim is that they do not occur
in the audited \texttt{*.tex} corpus under the stated repository path, not that a
priority search over the entire literature has been completed.
\end{abstract}

\medskip
\noindent\textbf{Status of results.}
Every result labelled theorem, proposition, lemma, or corollary below has a proof
in this report. The numerical tables are reproducible from
\path{verify_frontier_results.py}.  The denominator-free finite
\(q\)-binomial theorem is formalized in \path{FiniteQBinomialCore.lean}, while
\path{GeometricLagrange.lean} proves mass one, cancellation of the first
geometric modes, the first surviving moment, and module-valued Richardson
exactness.  The already compiled \path{GeometricQBinomialLagrange.lean}
identifies the individual weights with their Gaussian-binomial quotients under
injectivity of the finite geometric node set.  At post-snapshot source
checkpoint \path{1e761ed77583e9870ceeaeb2c74ac824094f0f3f},
\path{GeometricCompleteHomogeneous.lean} and
\path{GeometricResidualMoments.lean} additionally formalize the
denominator-free principal specialization and the complete positive-degree
Gaussian \(q\)-moment formula, including common scales and shifted blocks.
At the reconciled source checkpoint
\path{f77ff7c296b99050107948f5aa2a9488aec01d05},
the rational layer is implemented in
\begin{center}
\path{GeometricQBinomialLagrange.lean}\\
\path{GeometricRichardson.lean}\\
\path{GeometricLagrangeWeights.lean}\\
\path{GeometricLagrangeQBinomial.lean}\\
\path{GeometricLagrangeQMoments.lean}.
\end{center}
These modules prove the finite rational
\(0<q<1\) weight formulas, the \(W_m(q^r)\) and \(q\)-Pochhammer moment
presentations, all residual moments and their sign, and the exact finite
\(\ell^1\) norm, including dedicated \(q=1/4\) corollaries.  The
analytic sinc-product acceleration, transformed-error series, remainder and
bracketing bounds, spline limits, infinite condition-number limit, and
inverse-function consequences remain unformalized.

\tableofcontents

\section{Audit scope and the non-duplication boundary}

\subsection{How the corpus was used}

The source audit covered every LaTeX family in the repository documentation tree
\cite{repo}, including the current expository volumes, the formalization walkthrough, the
non-elementarity article, the mathematical glossary, the consolidated frontier
volume, the Thue--Morse formula atlas, all waves of the Fourier-decay discussion,
the paper-source subdirectories, and the historical archive. The archive index
was used to track renamed and byte-identical historical variants, so that an old
statement was not mistaken for a new one merely because it appeared under a
superseded filename. A source map is given in \cref{app:corpus}.

The audit revealed four especially important boundaries.

\begin{enumerate}[label=\textbf{B\arabic*.},leftmargin=12mm]
\item The main exposition and its supplements already contain the normalized
infinite sinc product, the box-convolution construction of \(\up\), its relation
to the Fabius function, dyadic identities, q-Pochhammer and q-binomial formulae,
and the signed Thue--Morse derivative comb.

\item The asymptotic documents already contain a corrected Lambert--\(W\)
endpoint scale, a periodic dyadic phase, and all-order small-argument expansions
for the Fabius function and its inverse. Those results concern
\(x\downarrow0\), not finite-product truncation \(n\to\infty\).

\item The Fourier-decay waves study \(\abs{P(\xi)}\) as
\(\abs{\xi}\to\infty\), including dyadic oscillation and exceptional
frequencies. They do not develop the exact analytic expansion of
\(P_{b,n}/P_b\) in the small geometric variable \(q^n\), nor the extrapolation
filter derived below.

\item The frontier volume already has q-Richardson ideas for discrete Fabius
spline rows and degree-dependent approximants. We therefore do not present a
generic Richardson recipe as new. The new object here is the \emph{Fourier-tail
truncation sequence} \(P_{b,n}\), for which the error coefficients, all q-moments,
stability constant, moment matching, and inverse-function consequence can all be
written exactly.
\end{enumerate}

The external search performed for this report also found general literature on
Richardson extrapolation and on powers of sinc expressed through Bell
polynomials, but no source matching the combined construction in
\cref{sec:tail,sec:filter,sec:physical}. This is supporting evidence only, not an
exhaustive priority claim.

\subsection{What is new in this report}

The principal contributions are the following.

\begin{enumerate}[label=\textbf{N\arabic*.},leftmargin=12mm]
\item A one-variable analytic tail function \(\Phi_b\) whose Taylor cumulants
are rational Bernoulli combinations and whose Taylor coefficients are complete
Bell polynomials.

\item A q-Lagrange operator
\(W_m(T_q)=\prod_{r=1}^m(T_q-q^r)/(1-q^r)\) that annihilates the first \(m\)
truncation modes exactly.

\item An all-order q-binomial error identity, not merely a leading-order
Richardson estimate.

\item A closed stability constant with a finite infinite-order limit, and a
successive-level error certificate requiring no knowledge of the infinite
product.

\item Exact matching of the first \(2m\) moments of the limiting Rvachev density
by a signed combination of finite box splines, for every \(n\).

\item A \(C^\infty\) superconvergence theorem with an explicit first derivative
profile, followed by an interior inverse-Fabius perturbation formula.

\item A kernel-independent extension showing that the same q-filter accelerates
any geometric product whose elementary factor has an even logarithmic Taylor
series.
\end{enumerate}

\section{Geometric sinc products and finite Rvachev splines}
\label{sec:setup}

Fix \(b>1\) and set
\[
 q=b^{-2}\in(0,1),\qquad \sinc z=\begin{cases}\sin z/z,&z\ne0,\\1,&z=0.\end{cases}
\]
Define
\begin{equation}
 P_b(\xi)=\prod_{\nu=1}^{\infty}\sinc(\xi/b^\nu),
 \qquad
 P_{b,n}(\xi)=\prod_{\nu=1}^{n}\sinc(\xi/b^\nu).
 \label{eq:products}
\end{equation}
The convergence is locally uniform in \(\C\), since
\(1-\sinc(\xi/b^\nu)=O(b^{-2\nu})\) on compact sets.

Let
\begin{equation}
 h_{b,\nu}(x)=\frac{b^\nu}{2}\,\mathbf 1_{[-b^{-\nu},b^{-\nu}]}(x),
 \qquad
 u_{b,n}=h_{b,1}*\cdots*h_{b,n}.
 \label{eq:boxes}
\end{equation}
With the Fourier convention
\[
 \widehat f(\xi)=\int_{\R}f(x)\e^{-\ii\xi x}\,dx,
 \qquad
 f(x)=\frac1{2\pi}\int_{\R}\widehat f(\xi)\e^{\ii\xi x}\,d\xi,
\]
we have \(\widehat h_{b,\nu}(\xi)=\sinc(\xi/b^\nu)\), hence
\begin{equation}
 \widehat{u_{b,n}}=P_{b,n}.
 \label{eq:finiteFT}
\end{equation}
The infinite convolution \(u_b=*_{\nu\geq1}h_{b,\nu}\) is a probability density
supported on \([-1/(b-1),1/(b-1)]\), and \(\widehat u_b=P_b\). The product
decays faster than every power, so \(u_b\in C_c^\infty(\R)\). At \(b=2\),
\(u_2\) is the normalized Rvachev up-function.

\subsection{The finite Thue--Morse derivative comb}

The bridge to the signed Thue--Morse sequence is already part of the repository
background, but the exact normalization is useful for what follows. Write
\(t_k=(-1)^{s_2(k)}\), where \(s_2(k)\) is the binary digit sum.

\begin{proposition}[distributional Thue--Morse comb]
For \(n\geq1\),
\begin{equation}
 \D^n u_{2,n}
 =2^{n(n-1)/2}\sum_{k=0}^{2^n-1}t_k\,
 \delta_{-1+(2k+1)2^{-n}}.
 \label{eq:TMcomb}
\end{equation}
\end{proposition}

\begin{proof}
For \(a>0\),
\[
 \D\left(\frac1{2a}\mathbf1_{[-a,a]}\right)
 =\frac1{2a}(\delta_{-a}-\delta_a).
\]
Differentiate every factor in \eqref{eq:boxes}. Choosing the right endpoint in
the \(j\)-th factor contributes one minus sign. If the choices are encoded by
the binary digits of \(k\), the coefficient is \((-1)^{s_2(k)}\), while the
location is
\[
 -\sum_{j=1}^n2^{-j}+2\sum_{j=1}^n\varepsilon_j2^{-j}
 =-1+2^{-n}+(2^{1-n})k.
\]
Finally \(\prod_{j=1}^n(2\cdot2^{-j})^{-1}=2^{n(n-1)/2}\).
\end{proof}

Thus every finite product in \eqref{eq:products} is simultaneously a Fourier
approximation and an integrated Thue--Morse object. The acceleration below acts
on both descriptions at once.

\section{The exact Bernoulli--Bell tail expansion}
\label{sec:tail}

\subsection{A universal tail function}

The key observation is that the two variables \(n\) and \(\xi\) enter the omitted
tail only through \(q^n\xi^2\).

\begin{definition}
For \(z\) near zero, define
\begin{equation}
 \Phi_b(z)=\prod_{\ell=1}^{\infty}
 \frac1{\sinc(b^{-\ell}\sqrt z)}.
 \label{eq:PhiDef}
\end{equation}
The expression is independent of the sign of \(\sqrt z\), and hence defines an
ordinary analytic function of \(z\).
\end{definition}

\begin{theorem}[exact tail factorization]
\label{thm:tailfactor}
The function \(\Phi_b\) is analytic in \(\abs z<\pi^2b^2\), meromorphic in
\(\C\), and has its nearest singularity to the origin at \(z=\pi^2b^2\). For
every \(n\geq0\),
\begin{equation}
 P_{b,n}(\xi)=P_b(\xi)\Phi_b(q^n\xi^2).
 \label{eq:exacttail}
\end{equation}
The identity is entire in \(\xi\); at common zeros it is interpreted by analytic
continuation. Moreover,
\begin{equation}
 \Phi_b(z)=\frac{\Phi_b(qz)}{\sinc(b^{-1}\sqrt z)}.
 \label{eq:PhiFunctional}
\end{equation}
\end{theorem}

\begin{proof}
Cancel the first \(n\) factors in \(P_b/P_{b,n}\) and write \(\nu=n+\ell\):
\[
 \frac{P_{b,n}(\xi)}{P_b(\xi)}
 =\prod_{\ell=1}^{\infty}
 \frac1{\sinc(\xi/b^{n+\ell})}
 =\Phi_b(q^n\xi^2).
\]
The zeros of \(\sinc(b^{-\ell}\sqrt z)\) occur at
\(z=\pi^2k^2b^{2\ell}\), \(k\in\mathbb Z\setminus\{0\}\). The nearest is
\(\pi^2b^2\). Splitting off the \(\ell=1\) factor and reindexing proves
\eqref{eq:PhiFunctional}.
\end{proof}

\subsection{Bernoulli cumulants}

The classical expansion
\begin{equation}
 \log\sinc w=-\sum_{k=1}^{\infty}
 \frac{\zeta(2k)}{k\pi^{2k}}w^{2k},\qquad \abs w<\pi,
 \label{eq:logsinc}
\end{equation}
turns the geometric tail into an explicitly summable cumulant series.

\begin{theorem}[Bernoulli cumulants]
\label{thm:cumulants}
For \(\abs z<\pi^2b^2\),
\begin{equation}
 \Phi_b(z)=\exp\!\left(\sum_{k=1}^{\infty}c_k(b)z^k\right),
 \label{eq:PhiCumulant}
\end{equation}
where
\begin{align}
 c_k(b)
 &=\frac{\zeta(2k)}{k\pi^{2k}(b^{2k}-1)}
 \label{eq:ckzeta}\\
 &=\frac{(-1)^{k+1}2^{2k-1}B_{2k}}
 {k(2k)!(b^{2k}-1)}.
 \label{eq:ckBernoulli}
\end{align}
In particular \(c_k(b)>0\) for every \(k\geq1\).
\end{theorem}

\begin{proof}
Insert \eqref{eq:logsinc} into \eqref{eq:PhiDef}, interchange absolutely
convergent sums, and use
\[
 \sum_{\ell=1}^{\infty}b^{-2k\ell}=\frac1{b^{2k}-1}.
\]
The Bernoulli form follows from Euler's formula for \(\zeta(2k)\). Positivity is
clear from \eqref{eq:ckzeta}.
\end{proof}

Combining \cref{thm:tailfactor,thm:cumulants} gives the exact truncation expansion
\begin{equation}
 \frac{P_{b,n}(\xi)}{P_b(\xi)}
 =\exp\!\left(\sum_{k\geq1}c_k(b)\xi^{2k}q^{kn}\right).
 \label{eq:tailn}
\end{equation}
This is an equality, not a merely formal asymptotic series, whenever
\(\abs\xi<\pi b^{n+1}\).

\subsection{Bell coefficients and their recurrence}

Write
\begin{equation}
 \Phi_b(z)=\sum_{r=0}^{\infty}\alpha_r(b)z^r,
 \qquad \alpha_0(b)=1.
 \label{eq:alphadef}
\end{equation}

\begin{proposition}[complete Bell-polynomial formula]
\label{prop:Bell}
Let \(Y_r\) denote the complete exponential Bell polynomial. Then
\begin{equation}
 \alpha_r(b)=\frac1{r!}
 Y_r\bigl(1!c_1(b),2!c_2(b),\ldots,r!c_r(b)\bigr),
 \label{eq:Bellalpha}
\end{equation}
and equivalently
\begin{equation}
 r\alpha_r(b)=\sum_{k=1}^r k c_k(b)\alpha_{r-k}(b).
 \label{eq:alpharec}
\end{equation}
Every \(\alpha_r(b)\) is strictly positive.
\end{proposition}

\begin{proof}
Equation \eqref{eq:Bellalpha} is the defining exponential generating identity for
complete Bell polynomials. Differentiating \eqref{eq:PhiCumulant} and comparing
coefficients gives \eqref{eq:alpharec}. Positivity follows inductively from
\(c_k(b)>0\).
\end{proof}

For the Rvachev base \(b=2\), the first cumulants are
\begin{align*}
 c_1&=\frac1{18},&c_2&=\frac1{2700},&c_3&=\frac1{178605},\\
 c_4&=\frac1{9639000},&c_5&=\frac1{478533825},&
 c_6&=\frac{691}{15688261338750}.
\end{align*}
The corresponding Bell coefficients are
\begin{align}
 \alpha_0&=1,&
 \alpha_1&=\frac1{18},&
 \alpha_2&=\frac{31}{16200},\\
 \alpha_3&=\frac{2347}{42865200},&
 \alpha_4&=\frac{1904369}{1311675120000},&
 \alpha_5&=\frac{3309760579}{88561680752160000}.
 \label{eq:alphaB2}
\end{align}
The exact rationals are useful because no numerical fitting is needed at any
order.

\section{q-Lagrange Richardson extrapolation}
\label{sec:filter}

\subsection{The filter polynomial}

Let \(T_q\) be the dilation operator \((T_qf)(z)=f(qz)\). Define
\begin{equation}
 W_m(t)=\prod_{r=1}^m\frac{t-q^r}{1-q^r},
 \qquad W_0(t)=1.
 \label{eq:Wm}
\end{equation}
Thus \(W_m(1)=1\) and \(W_m(q^r)=0\) for \(1\leq r\leq m\).
Expanding \(W_m(t)=\sum_{j=0}^m w_j^{(m)}t^j\) gives the following exact
weights.

\begin{theorem}[geometric Lagrange weights]
\label{thm:weights}
For \(0\leq j\leq m\),
\begin{equation}
 w_j^{(m)}
 =\prod_{\substack{0\leq\ell\leq m\\\ell\ne j}}
 \frac{-q^\ell}{q^j-q^\ell}
 =\frac{(-1)^{m-j}q^{(m-j)(m-j+1)/2}}
 {(q;q)_j(q;q)_{m-j}},
 \label{eq:weights}
\end{equation}
where \((q;q)_r=\prod_{s=1}^r(1-q^s)\). These are precisely the Lagrange
weights that recover the value at zero from data at \(1,q,\ldots,q^m\).
\end{theorem}

\begin{proof}
The first expression is the Lagrange cardinal basis evaluated at zero. Separating
factors with \(\ell<j\) and \(\ell>j\), then collecting powers of \(q\), gives the
second expression. Alternatively, the q-binomial theorem expands
\eqref{eq:Wm} directly.
\end{proof}

\paragraph{Lean boundary for the individual weights.}
The source-level closed quotient predates the new all-residual tranche.  Over a field, under
\(
 \operatorname{InjOn}(j\mapsto q^j,\{0,\ldots,m\}),
\)
the declaration
\begin{center}
\path{geometricLagrangeWeight_eq_gaussianBinomial_div}
\end{center}
in \path{GeometricQBinomialLagrange.lean} gives the numerator with the deliberate
reverse index \path{gaussianBinomial q m (m - j)}.  The displayed symmetric-
index notation in \eqref{eq:weights} is the equivalent paper convention, not a
claim that this checkpoint added a generic Gaussian-symmetry theorem.
\begin{samepage}
The same module also contains the following cleared-denominator identity:
\begin{center}
\path{finiteQPochhammerIn_mul_geometricLagrangeWeight_eq}.
\end{center}
\end{samepage}
In the present
real range \(0<q<1\), the required finite-node injectivity is automatic.  These
pointwise weight results should not be attributed to the later
principal-specialization module.

\begin{definition}[accelerated finite product]
For \(m,n\geq0\), define
\begin{equation}
 \mathcal R_{m,n}P_b(\xi)
 =\sum_{j=0}^m w_j^{(m)}P_{b,n+j}(\xi).
 \label{eq:Rdef}
\end{equation}
\end{definition}

With \(z=q^n\xi^2\), the tail factorization gives the compact operator identity
\begin{equation}
 \frac{\mathcal R_{m,n}P_b(\xi)}{P_b(\xi)}
 =\bigl(W_m(T_q)\Phi_b\bigr)(z).
 \label{eq:operatoridentity}
\end{equation}
The right side is well defined as an analytic continuation even at common zeros.

\subsection{A Romberg-style triangular algorithm}

The factorization of \(W_m\) gives a recursion that avoids an explicit weighted
sum.

\begin{algorithm}[q-Richardson table]
\label{alg:table}
Start with \(A_{0,n}=P_{b,n}(\xi)\). For \(m\geq1\), form
\begin{equation}
 A_{m,n}=\frac{A_{m-1,n+1}-q^mA_{m-1,n}}{1-q^m}.
 \label{eq:table}
\end{equation}
Then \(A_{m,n}=\mathcal R_{m,n}P_b(\xi)\).
\end{algorithm}

\begin{proof}
The polynomial recursion
\[
 W_m(t)=\frac{t-q^m}{1-q^m}W_{m-1}(t)
\]
becomes \eqref{eq:table} after applying it to the shift \(n\mapsto n+1\).
\end{proof}

For \(b=2\), \(q=1/4\), the first filters are
\begin{align}
 \mathcal R_{1,n}P&=\frac{4P_{n+1}-P_n}{3},
 \label{eq:m1}\\
 \mathcal R_{2,n}P&=\frac{P_n-20P_{n+1}+64P_{n+2}}{45},
 \label{eq:m2}\\
 \mathcal R_{3,n}P&=\frac{-P_n+84P_{n+1}-1344P_{n+2}+4096P_{n+3}}{2835}.
 \label{eq:m3}
\end{align}
The coefficients are exact rationals; the included program generates arbitrary
orders.

\subsection{The exact q-moment identity}

\begin{lemma}[all geometric moments]
\label{lem:qmoments}
For every integer \(r\geq0\), let
\(M_r^{(m)}=\sum_{j=0}^m w_j^{(m)}q^{jr}\). Then
\begin{equation}
 M_r^{(m)}=W_m(q^r).
 \label{eq:MrW}
\end{equation}
Consequently,
\begin{equation}
 M_0^{(m)}=1,
 \qquad M_r^{(m)}=0\quad(1\leq r\leq m),
 \label{eq:cancelmom}
\end{equation}
and for \(r\geq m+1\),
\begin{align}
 M_r^{(m)}
 &=q^{mr}\frac{(q^{1-r};q)_m}{(q;q)_m}
 \label{eq:Mrpoch}\\
 &=(-1)^m q^{\binom{m+1}{2}}\qbinom{r-1}{m}.
 \label{eq:Mrqbinom}
\end{align}
\end{lemma}

\begin{proof}
Equation \eqref{eq:MrW} follows from the coefficient definition of \(W_m\).
The roots give \eqref{eq:cancelmom}. For \(r>m\), evaluate the product in
\eqref{eq:Wm}:
\[
 W_m(q^r)=\prod_{\ell=1}^m\frac{q^r-q^\ell}{1-q^\ell}
 =(-1)^m q^{\binom{m+1}{2}}
 \prod_{\ell=1}^m\frac{1-q^{r-\ell}}{1-q^\ell},
\]
which is \eqref{eq:Mrqbinom}; \eqref{eq:Mrpoch} is the equivalent
q-Pochhammer form.
\end{proof}

\paragraph{Post-snapshot Lean status of the moment formula.}
The principal specialization
\[
 h_s(1,q,\ldots,q^m)=\qbinom{s+m}{m}
\]
is \path{completeHomogeneousEval_geometric}; it holds over every commutative
semiring, with no division or restriction on \(q\).  Its Lean right-hand side
is the recursive term \path{gaussianBinomial q (s + m) m}; in the present
real range this is the quotient denoted by \(\qbinom{s+m}{m}\).  More
generally, over a commutative ring,
\path{sum_weight_mul_geometric_pow_of_pos} proves that any row
\(a_0,\ldots,a_m\) satisfying
\[
 \sum_{j=0}^{m}a_j(q^j)^d=0^d\qquad(0\le d\le m)
\]
has, for every \(r>0\),
\[
 \sum_{j=0}^{m}a_j(q^j)^r
 =(-1)^m q^{\binom{m+1}{2}}\qbinom{r-1}{m}.
\]
Here and in the two formulas below, the literal Lean coefficient is the
recursive \path{gaussianBinomial}; the quotient notation is valid in this
report's real range \(0<q<1\).
The above-diagonal vanishing of the recursive Gaussian coefficient includes
all cancellations in \eqref{eq:cancelmom}; degree zero remains the separate
mass-one theorem.  For the Lagrange row, finite-node injectivity supplies the
low moments and \path{sum_geometricLagrangeWeight_mul_pow_of_pos} is exactly
the displayed formula.

The same checkpoint also records the scale rather than cancelling it.  If a
row is exact directly on \(c,cq,\ldots,cq^m\), meaning
\[
 \sum_{j=0}^{m}a_j(cq^j)^d=0^d\qquad(0\le d\le m),
\]
then \path{sum_weight_mul_scaled_geometric_pow_succ_add} gives, for \(s\ge0\),
\[
 \sum_{j=0}^{m}a_j(cq^j)^{m+1+s}
 =c^{m+1+s}(-1)^m q^{\binom{m+1}{2}}\qbinom{m+s}{m}.
\]
This hypothesis is imposed after scaling and remains meaningful when \(c\) is
zero or a zero divisor.  In the field-valued shifted block used here,
\path{sum_geometricLagrangeWeight_mul_shifted_pow_of_pos} specializes it to
\[
 \sum_{j=0}^{m}w_j^{(m)}(q^{n+j})^r
 =(-1)^m q^{nr+\binom{m+1}{2}}\qbinom{r-1}{m}
 \qquad(r>0).
\]
These are finite algebraic statements only; they do not establish the power-
series interchange or any analytic conclusion below.  All declarations in
this paragraph were checked at source checkpoint
\path{1e761ed77583e9870ceeaeb2c74ac824094f0f3f}.

\subsection{The complete error series}

\begin{theorem}[exact q-binomial error]
\label{thm:exacterror}
Let \(z=q^n\xi^2\), and suppose \(\abs z<\pi^2b^2\).  If
\(P_b(\xi)\ne0\), then
\begin{equation}
 \frac{\mathcal R_{m,n}P_b(\xi)}{P_b(\xi)}-1
 =(-1)^m q^{\binom{m+1}{2}}
 \sum_{r=m+1}^{\infty}
 \qbinom{r-1}{m}\alpha_r(b)z^r.
 \label{eq:exacterror}
\end{equation}
Without the nonvanishing hypothesis, the total statement is the multiplied
identity
\begin{equation}
 \mathcal R_{m,n}P_b(\xi)-P_b(\xi)
 =P_b(\xi)(-1)^m q^{\binom{m+1}{2}}
 \sum_{r=m+1}^{\infty}
 \qbinom{r-1}{m}\alpha_r(b)z^r.                         \label{eq:exacterror-multiplied}
\end{equation}
At a common zero of \(P_b\) and all the finite products, both sides of
\eqref{eq:exacterror-multiplied} vanish.  In particular, for fixed
\(m\) and fixed \(\xi\ne0\) with \(P_b(\xi)\ne0\), as \(n\to\infty\),
\begin{equation}
 \frac{\mathcal R_{m,n}P_b(\xi)}{P_b(\xi)}-1
 \sim (-1)^m q^{\binom{m+1}{2}}
 \alpha_{m+1}(b)\xi^{2m+2}q^{(m+1)n}.
 \label{eq:leadingerror}
\end{equation}
\end{theorem}

\begin{proof}
Apply the total tail identity \eqref{eq:exacttail} separately to each
\(P_{b,n+j}(\xi)\), multiply by \(w_j^{(m)}\), and sum.  This gives
\[
 \mathcal R_{m,n}P_b(\xi)
 =P_b(\xi)\sum_{j=0}^m w_j^{(m)}\Phi_b(q^jz)
\]
without division, including at zeros.  Insert
\(\Phi_b(z)=\sum_{r\geq0}\alpha_rz^r\); absolute convergence in the stated
disc permits the finite sum to pass through the series.  The coefficient of
\(z^r\) is \(\alpha_rW_m(q^r)\), so \cref{lem:qmoments} gives
\eqref{eq:exacterror-multiplied}.  Division by \(P_b(\xi)\) gives
\eqref{eq:exacterror} exactly on its stated locus.

Finally, analyticity at zero makes the residual series its first nonzero term
plus \(O(z^{m+2})\).  Since \(z=q^n\xi^2\), while
\(\alpha_{m+1}(b)>0\) by \cref{prop:Bell} and \(\xi\ne0\), that term is
nonzero and yields \eqref{eq:leadingerror}.
\end{proof}

This identity is substantially stronger than the statement that Richardson
extrapolation raises the order. Every residual mode and its exact q-binomial
multiplier are known.

\subsection{Divided differences and non-asymptotic sign bounds}

Let \(x_j=q^jz\). The quantity
\(\sum_jw_j^{(m)}\Phi_b(x_j)\) is the value at zero of the degree-\(m\)
Lagrange interpolant through the nodes \(x_0,\ldots,x_m\). The standard
interpolation remainder therefore yields a second exact form.

\begin{theorem}[divided-difference form]
\label{thm:divdiff}
For \(0\leq z<\pi^2b^2\),
\begin{equation}
 \bigl(W_m(T_q)\Phi_b\bigr)(z)-1
 =(-1)^m q^{\binom{m+1}{2}}z^{m+1}
 \Phi_b[0,q^mz,q^{m-1}z,\ldots,z],
 \label{eq:divdiff}
\end{equation}
where the bracket denotes a divided difference. Since \(\Phi_b\) is absolutely
monotone on this interval,
\begin{align}
 q^{\binom{m+1}{2}}\alpha_{m+1}(b)z^{m+1}
 &\leq
 \abs{\bigl(W_m(T_q)\Phi_b\bigr)(z)-1}
 \label{eq:lowerbound}\\
 &\leq
 \frac{q^{\binom{m+1}{2}}z^{m+1}}{(m+1)!}
 \Phi_b^{(m+1)}(z).
 \label{eq:upperbound}
\end{align}
The sign is exactly \((-1)^m\).
\end{theorem}

\begin{proof}
The Lagrange remainder at zero is
\[
 p_m(0)-\Phi_b(0)=(-1)^m\left(\prod_{j=0}^mx_j\right)
 \Phi_b[0,x_0,\ldots,x_m],
\]
and \(\prod_jx_j=q^{\binom{m+1}{2}}z^{m+1}\). Because every Taylor
coefficient \(\alpha_r\) is positive, all derivatives of \(\Phi_b\) are positive
and increasing on the stated interval. The Hermite--Genocchi formula bounds the
divided difference between
\(\Phi_b^{(m+1)}(0)/(m+1)!=\alpha_{m+1}\) and
\(\Phi_b^{(m+1)}(z)/(m+1)!\).
\end{proof}

\begin{corollary}[relative bracketing]
\label{cor:bracket}
For real \(\xi\) with \(\abs\xi<\pi b^{n+1}\) and \(P_b(\xi)\ne0\),
\begin{equation}
 \sgn\!\left(\frac{\mathcal R_{m,n}P_b(\xi)}{P_b(\xi)}-1\right)=(-1)^m.
 \label{eq:signlaw}
\end{equation}
Thus even and odd extrapolation orders bracket the exact product in relative
terms. If \(P_b(\xi)>0\), the bracketing is also an ordinary numerical ordering.
\end{corollary}

\subsection{Uniform stability of arbitrarily high order}

Richardson formulae are often numerically unattractive because their coefficient
norms grow rapidly. The geometric nodes here behave unusually well.

\begin{theorem}[exact Lebesgue constant]
\label{thm:stability}
The amplification factor of the filter is
\begin{equation}
 \Lambda_m(q)=\sum_{j=0}^m\abs{w_j^{(m)}}
 =\prod_{r=1}^m\frac{1+q^r}{1-q^r}
 =\frac{(-q;q)_m}{(q;q)_m}.
 \label{eq:Lebesgue}
\end{equation}
For fixed \(q\in(0,1)\),
\begin{equation}
 \Lambda_\infty(q)=\prod_{r=1}^{\infty}\frac{1+q^r}{1-q^r}<\infty.
 \label{eq:LebesgueInf}
\end{equation}
For \(q=1/4\),
\begin{equation}
 \Lambda_\infty(1/4)
 =1.969260353668269431947193696967\ldots.
 \label{eq:LambdaQuarter}
\end{equation}
\end{theorem}

\begin{proof}
The signs in \eqref{eq:weights} alternate, so
\[
 \sum_j\abs{w_j^{(m)}}=(-1)^mW_m(-1)
 =\prod_{r=1}^m\frac{1+q^r}{1-q^r}.
\]
The logarithm of the infinite product is
\(\sum_r[\log(1+q^r)-\log(1-q^r)]=O(\sum_rq^r)\), hence it converges.
\end{proof}

Consequently, absolute input errors bounded by \(\varepsilon\) produce an output
error no larger than \(\Lambda_\infty(q)\varepsilon\), uniformly in the
extrapolation order. This is a decisive difference from equally spaced
high-order extrapolation.

\subsection{A rigorous a posteriori certificate}

The positive-coefficient error series also compares consecutive truncation
levels.

\begin{theorem}[successive-level certificate]
\label{thm:certificate}
Fix \(m\) and a real \(\xi\). Assume
\(\abs\xi<\pi b^{n+1}\) and \(P_b(\xi)\ne0\). Let
\(R_n=\mathcal R_{m,n}P_b(\xi)\). Then
\begin{equation}
 \abs{R_{n+1}-R_n}
 \leq \abs{R_n-P_b(\xi)}
 \leq \frac{\abs{R_{n+1}-R_n}}{1-q^{m+1}}.
 \label{eq:certificate}
\end{equation}
Furthermore,
\begin{equation}
 P_b(\xi)-R_n
 \sim \frac{R_{n+1}-R_n}{1-q^{m+1}}.
 \label{eq:asympest}
\end{equation}
\end{theorem}

\begin{proof}
By \eqref{eq:exacterror}, the relative error is \((-1)^mA_n\), where
\[
 A_n=q^{\binom{m+1}{2}}
 \sum_{r\geq m+1}\qbinom{r-1}{m}\alpha_r\xi^{2r}q^{rn}>0.
\]
Every term is multiplied by \(q^r\) when \(n\) increases by one. Hence
\(0<A_{n+1}\leq q^{m+1}A_n\). Therefore
\[
 (1-q^{m+1})A_n\leq A_n-A_{n+1}\leq A_n.
\]
Multiplying by \(\abs{P_b(\xi)}\) proves \eqref{eq:certificate}. The leading
term has ratio \(q^{m+1}\), which proves \eqref{eq:asympest}.
\end{proof}

This certificate is practical: it estimates the unknown infinite product from
two consecutive accelerated values and uses only the known number
\(1-q^{m+1}\).

\section{Physical-space consequences for Rvachev splines}
\label{sec:physical}

Define the signed extrapolated spline
\begin{equation}
 U_{b;m,n}(x)=\sum_{j=0}^m w_j^{(m)}u_{b,n+j}(x).
 \label{eq:Udef}
\end{equation}
It has total mass one and support in \([-1/(b-1),1/(b-1)]\). It need not be
nonnegative, but
\begin{equation}
 \norm{U_{b;m,n}}_{L^1}\leq\Lambda_m(q)<\Lambda_\infty(q).
 \label{eq:L1stable}
\end{equation}
Its Fourier transform is \(\mathcal R_{m,n}P_b\).

\subsection{All-order inverse-differential expansion}

\begin{theorem}[all-order spline expansion]
\label{thm:CinftyExpansion}
For every pair of nonnegative integers \(N,s\),
\begin{equation}
 u_{b,n}
 =\sum_{r=0}^{N}(-1)^r\alpha_r(b)q^{rn}u_b^{(2r)}
 +O_{C^s}\!\left(q^{(N+1)n}\right)
 \qquad(n\to\infty).
 \label{eq:physicalExpansion}
\end{equation}
Equivalently, as an all-order asymptotic operator identity,
\begin{equation}
 u_{b,n}\sim
 \exp\!\left(\sum_{k\geq1}(-1)^kc_k(b)q^{kn}\D^{2k}\right)u_b.
 \label{eq:operatorPhysical}
\end{equation}
\end{theorem}

\begin{proof}
Fourier transformation turns \((-1)^r\D^{2r}u_b\) into
\(\xi^{2r}P_b(\xi)\), so the displayed finite sum is exactly the first
\(N+1\) terms of \eqref{eq:tailn}.

For completeness, split the inverse Fourier integral at \(\abs\xi=b^n\). On the
low-frequency region, \(z=q^n\xi^2\leq1\), and analyticity of \(\Phi_b\) gives a
uniform Taylor remainder bounded by
\(Cq^{(N+1)n}\abs\xi^{2N+2}\abs{P_b(\xi)}\). This remains integrable after
multiplication by \(\abs\xi^s\), because \(P_b\) is rapidly decreasing.

On \(\abs\xi\geq b^n\), the first \(n\) factors satisfy
\[
 \abs{P_{b,n+j}(\xi)}\leq
 b^{n(n+1)/2}\abs\xi^{-n},\qquad j\geq0,
\]
while the remaining factors have modulus at most one. The corresponding weighted
Fourier tail is \(O(b^{-n^2/2+O(n)})\), hence smaller than every fixed geometric
rate. The same bound applies to \(P_b\). Fourier inversion proves the assertion
in \(C^s\).
\end{proof}

\subsection{Superconvergence of the extrapolated spline}

\begin{theorem}[q-binomial \(C^\infty\) expansion]
\label{thm:Uexpansion}
For fixed \(m,s\) and every \(N\geq m+1\),
\begin{align}
 U_{b;m,n}-u_b
 &=q^{\binom{m+1}{2}}
 \sum_{r=m+1}^{N}(-1)^{m+r}
 \qbinom{r-1}{m}\alpha_r(b)q^{rn}u_b^{(2r)}
 \notag\\
 &\quad+O_{C^s}\!\left(q^{(N+1)n}\right).
 \label{eq:Uallorder}
\end{align}
In particular,
\begin{equation}
 \norm{U_{b;m,n}-u_b}_{C^s}
 \leq C_{b,m,s}q^{(m+1)n},
 \label{eq:Csuper}
\end{equation}
and the normalized error has the explicit limit
\begin{equation}
 q^{-(m+1)n}\bigl(U_{b;m,n}-u_b\bigr)
 \longrightarrow
 -q^{\binom{m+1}{2}}\alpha_{m+1}(b)u_b^{(2m+2)}
 \quad\text{in }C^s(\R).
 \label{eq:errorprofile}
\end{equation}
\end{theorem}

\begin{proof}
Apply inverse Fourier transformation to \eqref{eq:exacterror}, using
\(\mathcal F^{-1}[\xi^{2r}P_b]=(-1)^ru_b^{(2r)}\). The remainder estimate is
the same low/high-frequency argument as in \cref{thm:CinftyExpansion}. At
\(r=m+1\), the sign is \((-1)^{m+m+1}=-1\), which gives
\eqref{eq:errorprofile}.
\end{proof}

For \(b=2\), each additional extrapolation level changes the basic convergence
factor from \(4^{-n}\) to \(4^{-(m+1)n}\), while the filter norm remains below
\(1.970\).

\subsection{Exact moment matching}

The cancellation is stronger than fast convergence: a finite set of moments is
already exact.

\begin{theorem}[exact moment matching]
\label{thm:moments}
For every \(m,n\geq0\) and every integer \(0\leq\ell\leq2m\),
\begin{equation}
 \int_{\R}x^\ell U_{b;m,n}(x)\,dx
 =\int_{\R}x^\ell u_b(x)\,dx.
 \label{eq:momentmatch}
\end{equation}
All odd moments vanish by symmetry. The first missed even moment is known
exactly:
\begin{equation}
 \int_{\R}x^{2m+2}\bigl(U_{b;m,n}(x)-u_b(x)\bigr)\,dx
 =-(2m+2)!\,q^{\binom{m+1}{2}}
 \alpha_{m+1}(b)q^{(m+1)n}.
 \label{eq:firstmomenterror}
\end{equation}
\end{theorem}

\begin{proof}
By \eqref{eq:exacterror},
\(\widehat U_{b;m,n}-\widehat u_b\) has a zero of order at least \(2m+2\) at
\(\xi=0\). Since
\[
 \widehat f^{(\ell)}(0)=(-\ii)^\ell\int x^\ell f(x)\,dx,
\]
all moments through degree \(2m+1\) agree. Symmetry handles every odd degree.
The coefficient of \(\xi^{2m+2}\) in the Fourier error is
\((-1)^mq^{\binom{m+1}{2}}\alpha_{m+1}q^{(m+1)n}\). Converting the derivative
at zero back to the \((2m+2)\)-nd moment contributes \((-1)^{m+1}\), leaving the
negative sign in \eqref{eq:firstmomenterror}.
\end{proof}

\begin{remark}[quadrature interpretation]
For any polynomial \(p\) of degree at most \(2m\),
\[
 \int p(x)U_{b;m,n}(x)\,dx=\int p(x)u_b(x)\,dx.
\]
Thus \(U_{b;m,n}\) is a signed, compactly supported, uniformly stable density
surrogate with exact polynomial expectations through degree \(2m\).
\end{remark}

\section{Fabius and inverse-Fabius consequences}
\label{sec:fabius}

Under the normalization used throughout the repository, the Fabius function on
\([0,1]\) is the left half of the translated up-function:
\begin{equation}
 F(x)=u_2(x-1),\qquad 0\leq x\leq1.
 \label{eq:Fshift}
\end{equation}
Define the accelerated finite-spline approximation
\begin{equation}
 F_{m,n}(x)=U_{2;m,n}(x-1),\qquad q=\frac14.
 \label{eq:Fapprox}
\end{equation}

\begin{corollary}[accelerated Fabius approximation]
\label{cor:Fconv}
For every \(m,s\geq0\),
\begin{equation}
 \norm{F_{m,n}-F}_{C^s([0,1])}
 =O\!\left(4^{-(m+1)n}\right),
 \label{eq:Fconv}
\end{equation}
and
\begin{equation}
 4^{(m+1)n}(F_{m,n}-F)
 \longrightarrow
 -4^{-\binom{m+1}{2}}\alpha_{m+1}(2)F^{(2m+2)}
 \quad\text{in }C^s([0,1]).
 \label{eq:Fprofile}
\end{equation}
\end{corollary}

The global statement is possible because all derivatives of the compactly
supported up-function vanish at the endpoints. Inversion is more delicate:
\(F'(x)\) tends to zero at the ends, and this is precisely where the repository's
Lambert--\(W\) asymptotics become relevant.

\begin{theorem}[interior inverse-Fabius superconvergence]
\label{thm:inverse}
Let \(K=[a,b]\Subset(0,1)\), and put \(J=F(K)\). Let \(G=F^{-1}\) on \(J\).
For sufficiently large \(n\), the restriction of \(F_{m,n}\) to a neighborhood
of \(K\) is strictly increasing and has an inverse branch \(G_{m,n}\) on \(J\).
Then
\begin{equation}
 \norm{G_{m,n}-G}_{C^s(J)}
 =O\!\left(4^{-(m+1)n}\right)
 \label{eq:Gconv}
\end{equation}
for every fixed \(s\), and uniformly for \(y\in J\),
\begin{equation}
 4^{(m+1)n}\bigl(G_{m,n}(y)-G(y)\bigr)
 \longrightarrow
 4^{-\binom{m+1}{2}}\alpha_{m+1}(2)
 \frac{F^{(2m+2)}(G(y))}{F'(G(y))}.
 \label{eq:Gprofile}
\end{equation}
\end{theorem}

\begin{proof}
On a slightly larger compact interval \(K'\Subset(0,1)\), continuity and strict
monotonicity give \(F'\geq c>0\). The \(C^1\) convergence in
\cref{cor:Fconv} implies \(F_{m,n}'\geq c/2\) for large \(n\), so the inverse
branch exists. Standard differentiation of inverse functions gives
\eqref{eq:Gconv}.

For the leading term, write \(F_{m,n}=F+\varepsilon_n\), set \(x=G(y)\), and
solve
\[
 0=F(x+\delta_n)-F(x)+\varepsilon_n(x+\delta_n)
 =F'(x)\delta_n+\varepsilon_n(x)+o(4^{-(m+1)n}).
\]
Equation \eqref{eq:Fprofile} gives
\(\varepsilon_n(x)\sim
 -4^{-\binom{m+1}{2}}\alpha_{m+1}4^{-(m+1)n}F^{(2m+2)}(x)\),
which proves \eqref{eq:Gprofile}.
\end{proof}

\subsection{Relation to \texorpdfstring{Lambert--\(W\)}{Lambert-W} endpoint asymptotics}

The inverse theorem is intentionally restricted to compact interior intervals.
Near \(y=0\), the factor \(1/F'(G(y))\) becomes large and fixed-order inversion
is no longer uniform. The repository's endpoint theory expresses the dominant
scale of \(G(y)\) using a Lambert--\(W\) variable together with a periodic dyadic
phase. The two results therefore occupy complementary limits:
\[
 \begin{array}{c|c|c}
 \text{regime}&\text{small parameter}&\text{governing mechanism}\\ \hline
 y\downarrow0 & 1/\log(1/y) & \text{Lambert--}W\text{ and dyadic phase}\\
 n\to\infty,\ y\in J\Subset(0,1)&4^{-n}&\text{Bernoulli--Bell tail and q-filter}.
 \end{array}
\]
A genuinely new two-parameter problem is to choose \(n=n(y)\) so that
\eqref{eq:Gprofile} remains uniform while \(y\downarrow0\). The exact factor
\(F^{(2m+2)}(G(y))/F'(G(y))\), combined with the repository's all-order endpoint
expansion, reduces that problem to a concrete matching calculation rather than a
vague stability question; see \cref{sec:open}.

\section{Numerical verification}
\label{sec:numerics}

All computations in this section used 90-decimal-digit arithmetic. The infinite
product was continued until a geometric bound for the omitted logarithmic tail
fell below \(10^{-70}\). The explicit weighted formula and the triangular table
\eqref{eq:table} were evaluated independently and agreed to the working
precision. Symbolic checks verified \eqref{eq:Mrqbinom} for a range of orders as
an identity in the symbol \(q\).

\subsection{Absolute errors}

For \(b=2\), the reference values are
\begin{align*}
 P(1)&=0.9456037838988178752649922474364995\ldots,\\
 P(3)&=0.5857191765350684829125292624510928\ldots,\\
 P(10)&=-0.03192757801277845828274435480432725\ldots.
\end{align*}

\begin{table}[ht]
\centering
\caption{Absolute errors at \(\xi=1\).}
\label{tab:error1}
\begin{tabular}{ccccc}
\toprule
\(n\)&\(m=0\)&\(m=1\)&\(m=2\)&\(m=3\)\\
\midrule
3 & \(8.2128\!\times10^{-4}\) & \(1.1050\!\times10^{-7}\) & \(3.0877\!\times10^{-12}\) & \(1.9989\!\times10^{-17}\)\\
4 & \(2.0524\!\times10^{-4}\) & \(6.9036\!\times10^{-9}\) & \(4.8225\!\times10^{-14}\) & \(7.8050\!\times10^{-20}\)\\
5 & \(5.1304\!\times10^{-5}\) & \(4.3143\!\times10^{-10}\) & \(7.5345\!\times10^{-16}\) & \(3.0485\!\times10^{-22}\)\\
6 & \(1.2826\!\times10^{-5}\) & \(2.6964\!\times10^{-11}\) & \(1.1772\!\times10^{-17}\) & \(1.1908\!\times10^{-24}\)\\
\bottomrule
\end{tabular}
\end{table}

\begin{table}[ht]
\centering
\caption{Absolute errors at \(\xi=3\).}
\label{tab:error3}
\begin{tabular}{ccccc}
\toprule
\(n\)&\(m=0\)&\(m=1\)&\(m=2\)&\(m=3\)\\
\midrule
3 & \(4.5982\!\times10^{-3}\) & \(5.5691\!\times10^{-6}\) & \(1.4003\!\times10^{-9}\) & \(8.1582\!\times10^{-14}\)\\
4 & \(1.1454\!\times10^{-3}\) & \(3.4676\!\times10^{-7}\) & \(2.1800\!\times10^{-11}\) & \(3.1753\!\times10^{-16}\)\\
5 & \(2.8608\!\times10^{-4}\) & \(2.1652\!\times10^{-8}\) & \(3.4031\!\times10^{-13}\) & \(1.2392\!\times10^{-18}\)\\
6 & \(7.1504\!\times10^{-5}\) & \(1.3529\!\times10^{-9}\) & \(5.3162\!\times10^{-15}\) & \(4.8397\!\times10^{-21}\)\\
\bottomrule
\end{tabular}
\end{table}

\begin{table}[ht]
\centering
\caption{Absolute errors at \(\xi=10\), beyond the first sign change of the
infinite product but still inside the tail disk for the displayed \(n\).}
\label{tab:error10}
\begin{tabular}{ccccc}
\toprule
\(n\)&\(m=0\)&\(m=1\)&\(m=2\)&\(m=3\)\\
\midrule
3 & \(2.9276\!\times10^{-3}\) & \(3.9468\!\times10^{-5}\) & \(1.1012\!\times10^{-7}\) & \(7.1217\!\times10^{-11}\)\\
4 & \(7.0230\!\times10^{-4}\) & \(2.3635\!\times10^{-6}\) & \(1.6504\!\times10^{-9}\) & \(2.6705\!\times10^{-13}\)\\
5 & \(1.7380\!\times10^{-4}\) & \(1.4617\!\times10^{-7}\) & \(2.5525\!\times10^{-11}\) & \(1.0327\!\times10^{-15}\)\\
6 & \(4.3341\!\times10^{-5}\) & \(9.1120\!\times10^{-9}\) & \(3.9782\!\times10^{-13}\) & \(4.0240\!\times10^{-18}\)\\
\bottomrule
\end{tabular}
\end{table}

The observed reduction on increasing \(n\) is approximately
\(4^{-(m+1)}\), exactly as predicted.

\subsection{Leading constants}

At \(\xi=3\), divide the signed error by \(q^{(m+1)n}\). The predicted limits
from \eqref{eq:leadingerror} are
\begin{center}
\begin{tabular}{ccc}
\toprule
\(m\)&predicted limit&computed at \(n=11\)\\
\midrule
0 & \(0.29285958826753424146\) & \(0.29285960991271731281\)\\
1 & \(-0.022696618090733903713\) & \(-0.022696619832602340475\)\\
2 & \(0.000365296265765254499\) & \(0.000365296293045216575\)\\
3 & \(-1.36214818263598388\times10^{-6}\) & \(-1.36214828256045162\times10^{-6}\)\\
\bottomrule
\end{tabular}
\end{center}
The alternating signs and limiting constants are resolved without fitting any
free parameter.

\subsection{A posteriori estimates}

For example, at \(\xi=3,n=4,m=3\), the true absolute error is
\(3.17531383672\times10^{-16}\). The computable difference between successive
accelerated values is \(3.16292144780\times10^{-16}\), and
\cref{thm:certificate} gives the upper bound
\(3.17532506132\times10^{-16}\). The certificate therefore encloses the true
error in an interval of relative width about \(4\times10^{-6}\), using no
reference value of \(P(3)\).

\section{A universal geometric-product theorem}
\label{sec:universal}

The q-filter is not tied to the sine function. Its universality clarifies which
parts of the construction are special to Rvachev's function.

\begin{theorem}[universal analytic-factor acceleration]
\label{thm:universal}
Let \(\varphi\) be an even analytic function near zero with \(\varphi(0)=1\),
nonzero in a neighborhood of zero, and
\begin{equation}
 \log\varphi(w)=-\sum_{k=1}^{\infty}a_kw^{2k}.
 \label{eq:generalphi}
\end{equation}
For \(b>1\), define
\[
 Q_b(\xi)=\prod_{\nu\geq1}\varphi(\xi/b^\nu),
 \qquad Q_{b,n}(\xi)=\prod_{\nu=1}^n\varphi(\xi/b^\nu).
\]
Then locally at \(q^n\xi^2=0\),
\begin{equation}
 \frac{Q_{b,n}(\xi)}{Q_b(\xi)}
 =\exp\!\left(\sum_{k\geq1}
 \frac{a_k}{b^{2k}-1}\xi^{2k}q^{kn}\right).
 \label{eq:universalTail}
\end{equation}
Let its Taylor coefficients in \(z=q^n\xi^2\) be \(\beta_r\). The same weights
\eqref{eq:weights} satisfy
\begin{equation}
 \frac{\sum_{j=0}^mw_j^{(m)}Q_{b,n+j}(\xi)}{Q_b(\xi)}-1
 =(-1)^mq^{\binom{m+1}{2}}
 \sum_{r\geq m+1}\qbinom{r-1}{m}\beta_r z^r.
 \label{eq:universalError}
\end{equation}
If all \(a_k\geq0\), then all \(\beta_r\geq0\), and the sign law, divided-
difference bounds, and a posteriori certificate remain valid.
\end{theorem}

\begin{proof}
The proof of \cref{thm:cumulants} uses only the geometric sum
\(\sum_{\ell\geq1}b^{-2k\ell}=1/(b^{2k}-1)\). The q-filter calculation then
uses only the Taylor expansion and \cref{lem:qmoments}.
\end{proof}

Thus the sinc-specific information is concentrated in the Bernoulli values of
\(a_k=\zeta(2k)/(k\pi^{2k})\). The q-binomial extrapolation, its stability, and
its exact mode cancellation belong to the geometry of the scales rather than to
the elementary factor.

\section{Lean formalization roadmap}
\label{sec:lean}

The analytic proof naturally separates into formalization-sized modules.

\begin{enumerate}[label=\arabic*.,leftmargin=9mm]
\item \textbf{Finite q-algebra (formalized).}  The finite \(q\)-Pochhammer
and Gaussian-binomial core is in \path{FiniteQBinomialCore.lean}; the closed
weights are in \path{GeometricQBinomialLagrange.lean}; and the principal
specialization, all positive moments, scaled exact-row formula, and shifted
Lagrange specialization are in \path{GeometricCompleteHomogeneous.lean} and
\path{GeometricResidualMoments.lean}.  The report-facing rational geometric
modules cited above, at checkpoint
\path{f77ff7c296b99050107948f5aa2a9488aec01d05}, additionally identify the
weight polynomial with \(W_m\), prove the \(q\)-Pochhammer moment presentation,
the sign of every residual moment, the exact finite norm, and dedicated
quarter-base corollaries.  In particular the
cancellations \(1\leq r\leq m\) and \eqref{eq:Mrqbinom} have exact Lean
counterparts.

\item \textbf{Formal power series.} Work first in a characteristic-zero
commutative ring. From \(\Phi=\exp(\sum c_kz^k)\), derive the Bell recurrence
\eqref{eq:alpharec}; then prove the coefficientwise filtered identity
\eqref{eq:exacterror}. This avoids analytic convergence in the first pass.

\item \textbf{Analytic sinc logarithm.} Establish \eqref{eq:logsinc} in a disk,
sum the geometric scales, and identify the analytic power series with the actual
product. Existing repository lemmas about the sinc product should reduce the
amount of new infrastructure.

\item \textbf{Interpolation and positivity (partially formalized).}  The
finite Lagrange weights and exact residual sign are proved.  The
divided-difference representation of the analytic transformed error and the
absolute-monotonicity argument for its full positive coefficient series remain
frontier work.

\item \textbf{Fourier/measure layer.} Define the uniform box measures and their
finite convolutions. Moment matching follows from derivatives of the
characteristic function at zero and may be formalized before the full
\(C^\infty\) theorem.

\item \textbf{Asymptotic layer.} Prove low-frequency Taylor bounds and the
high-frequency superexponential estimate. This is the technically heaviest step,
but it is independent of q-binomial algebra.

\item \textbf{Inverse-function layer.} Apply a quantitative inverse function
theorem on compact intervals where \(F'\) has a positive lower bound, then derive
\eqref{eq:Gprofile} from the first-order perturbation identity.
\end{enumerate}

The accompanying Python file supplies exact test vectors for the algebraic
modules and high-precision regression tests for the analytic statements.

\section{Further frontier problems}
\label{sec:open}

\subsection{A matched Lambert--Richardson endpoint theory}

Let \(G=F^{-1}\) and let \(G_{m,n}\) be the accelerated inverse branch. Determine
the largest endpoint region \(0<y\leq y_n\) on which
\[
 G_{m,n}(y)-G(y)
 \sim 4^{-\binom{m+1}{2}}\alpha_{m+1}(2)4^{-(m+1)n}
 \frac{F^{(2m+2)}(G(y))}{F'(G(y))}
\]
remains uniform. Substitution of the repository's Lambert--\(W\) expansion for
\(G(y)\) should identify a sharp transition curve \(y_n\), together with a
dyadically periodic modulation. This would connect two asymptotic theories that
are presently separate.

\subsection{Optimal order in finite precision}

Although \(\Lambda_m(1/4)<1.970\), the finite products themselves may suffer
relative loss near zeros. A complete finite-precision theory should choose
\(m\) and \(n\) adaptively from: (i) the certificate \eqref{eq:certificate};
(ii) the distance to the shared zero set; and (iii) the working precision.
Because the filter norm is uniformly bounded, one expects the optimal order to be
governed mainly by product evaluation rather than extrapolation instability.

\subsection{Positivity restoration}

The signed density \(U_{b;m,n}\) has exact low moments but need not be positive.
Is there a nonlinear, positivity-preserving correction that retains
\(q^{(m+1)n}\) convergence and the moment constraints? One possibility is a
small entropy projection onto the convex set of nonnegative densities with the
same moments. The exact first missed moment \eqref{eq:firstmomenterror} provides
a natural scale for such a correction.

\subsection{Dyadic rational evaluation from accelerated splines}

At a dyadic point, each finite spline \(u_{2,n+j}\) is exactly computable from
integrated Thue--Morse data and q-binomial formulae already present in the
repository. Substituting those values into \eqref{eq:Rdef} should yield explicit
rational or q-rational approximants to \(F(x)\) with a certified
\(4^{-(m+1)n}\) remainder. This would produce a new bridge between exact dyadic
arithmetic and the analytic sinc tail.

\subsection{Complex domains and pole-aware extrapolation}

The nearest tail pole at \(z=\pi^2b^2\) controls the Taylor expansion. For
complex \(\xi\), one can conformally map the slit \(z\)-plane before applying a
geometric interpolation filter. A pole-aware variant may retain exact q-mode
cancellation while extending useful acceleration beyond the first omitted sinc
zero.

\subsection{Other refinable and compactly supported functions}

Apply \cref{thm:universal} to elementary factors arising from compactly
supported refinable distributions, Bernoulli convolutions, or symmetric
B-splines. The main classification question is: which elementary factors have
nonnegative logarithmic cumulants, and therefore inherit the exact sign law and
successive-level certificate?

\section{Conclusion}

The truncation of the Rvachev sinc product carries more structure than a generic
convergent product. Its omitted tail is an analytic function of the single
geometric variable \(z=4^{-n}\xi^2\). Bernoulli numbers describe the logarithm
of that function, Bell polynomials describe the function itself, and Lagrange
interpolation on the nodes \(1,q,\ldots,q^m\) produces a q-binomial Richardson
filter with a complete, sign-definite error series.

The most notable consequences are exact rather than heuristic: the q-moments of
the weights are known in closed form; the filter norm stays bounded at arbitrarily
high order; consecutive accelerated levels certify the error; and the associated
signed box splines match the first \(2m\) moments of the limiting up-density for
every \(n\). Fourier inversion converts the same algebra into
\(C^\infty\)-superconvergence with an explicit derivative profile. Translation
and local inversion then give a new interior approximation theory for the Fabius
function and its inverse, complementary to the repository's Lambert--\(W\)
endpoint asymptotics.

These results create concrete next steps: continue the Lean formalization from
the compiled finite \(q\)-algebra into the formal-series and analytic error
layers, build a dyadic exact-evaluation implementation, and develop a matched
two-parameter theory linking truncation acceleration to the periodic Lambert
phase near the endpoint.

\appendix

\section{Exact coefficients and filters at \texorpdfstring{\(b=2\)}{b=2}}
\label{app:coeff}

The next two Bell coefficients, useful for regression tests, are
\begin{align*}
 \alpha_6&=\frac{1884491755011469}
 {1980124051433319792000000},\\
 \alpha_7&=\frac{151616054172597829}
 {6278781742186853835936000000}.
\end{align*}

The order-four filter is
\begin{align*}
 \mathcal R_{4,n}P
 &=\frac1{722925}P_n-\frac4{8505}P_{n+1}
 +\frac{64}{2025}P_{n+2}\\
 &\qquad-\frac{4096}{8505}P_{n+3}
 +\frac{1048576}{722925}P_{n+4}.
\end{align*}
The order-five weights are
\[
 \left(
 -\frac1{739552275},\frac4{2168775},-\frac{64}{127575},
 \frac{4096}{127575},-\frac{1048576}{2168775},
 \frac{1073741824}{739552275}
 \right).
\]
Despite the large integer numerators, the sum of absolute values is below
\(1.968\) at order five.

\section{Corpus audit map}
\label{app:corpus}

This appendix records the source families used to establish the non-duplication
boundary. Historical variants were read for content rather than treated as
independent discoveries.

\begingroup
\small
\raggedright
\subsection*{Current top-level expositions}
\begin{itemize}[leftmargin=8mm]
\item \path{Fabius_Function_and_Rvachev_Up/Fabius_Function_and_Rvachev_Up.tex}
\item \path{FabiusFunction_Mathematical_Glossary/FabiusFunction_Mathematical_Glossary.tex}
\item \path{Non_Elementarity_of_the_Fabius_Function/Non_Elementarity_of_the_Fabius_Function.tex}
\item \path{fabius_lean_walkthrough/fabius_lean_walkthrough.tex}
\item \path{non-formalized-research-frontiers/non-formalized-research-frontiers.tex}
\end{itemize}

\subsection*{Current frontier drafts}
\begin{itemize}[leftmargin=8mm]
\item \path{drafts/Thue_Morse_Formula_Atlas/Thue_Morse_Formula_Atlas.tex}
\item LaTeX source in \path{Rvachev_Up_Fourier_Decay-2}
\item LaTeX source in \path{Rvachev_Up_Fourier_Decay-3}
\item LaTeX source in \path{Rvachev_Up_Fourier_Decay_Comparative_Audit}
\item LaTeX source in \path{Rvachev_Up_Fourier_Decay_Gentle_Guide}
\item LaTeX source in \path{Rvachev_Up_Fourier_Decay_Second_Wave_Audit}
\item LaTeX source in \path{asymptotic-decay-rate-of-an-infinite-product-of-sinc-functions}
\item LaTeX source in \path{rvachev_up_fourier_decay-1}
\item LaTeX sources in \path{rvachev_up_fourier_decay-4} through
\path{rvachev_up_fourier_decay-8}
\end{itemize}

\subsection*{Paper-source directories}
\begin{itemize}[leftmargin=8mm]
\item \path{AIF_2017__67_2_637_0.tex} and its corrected source
\item \path{Lacunary_products.tex} (arXiv:1502.06738v2)
\item \path{09-Function.tex} (arXiv:1702.05442v1)
\item \path{157-Arithmetic-v3.tex} (arXiv:1702.06487v3)
\item \path{document.tex} (arXiv:2211.13570v1)
\item \path{ubiq15-corrected.tex}
\end{itemize}

\subsection*{Historical archive families}
\begin{itemize}[leftmargin=8mm]
\item \texttt{Fabius Asymptotic.tex}
\item \texttt{K-fold summation over signed Thue-Morse.tex}
\item \path{Non_Elementarity_of_the_Fabius_Function.tex}
\item \path{Small_Argument_Asymptotics.tex}
\item \path{fabius-inverse-dyadic-closed-form.tex}
\item Historical primary exposition variants under
\path{Fabius_Function_and_Rvachev_Up}, \path{Fabius_and_Rvachev_Up-2},
\path{Fabius_and_Rvachev_up}, and \path{Fabius_function_and_up_function}
\item \path{Fabius_Function_Missing_Parts} variants
\item \path{Fabius_q_Special_Functions} and \path{fabius-q-special-functions}
\item \path{Fabius_Function_Lean_Walkthrough} and
\path{fabius_lean_walkthrough} variants
\item \path{FabiusFunction_Mathematical_Glossary} variants
\item \path{Fabius_Dyadic_Asymptotic_Bridge}
\item \path{Fabius_Dyadic_Formulae_and_Alternative_Representations}
\item \path{Fabius_Dyadic_Formulae_to_Asymptotics}
\item \path{Fabius_Dyadic_q_Connections}
\item \path{Fabius_Thue_Morse_Convergence_Rate}
\item \path{Fabius_Thue_Morse_Convergence_Rates}
\item \path{Primary_Exposition_Gap_Register}
\item \path{Repeated_Integration_and_Rvachev_Up}
\item \path{Rvachev_Up_from_Repeated_Integration}
\item Earlier consolidated \path{non-formalized-research-frontiers} source
\item \path{Thue_Morse_Formula_Atlas-1}, \path{Thue_Morse_Formula_Atlas-2}, and
\path{Thue_Morse_Formula_Atlas-3}
\end{itemize}
\endgroup

The current consolidated frontier source explicitly records the absorption of
several former notebooks and their checksums. That provenance was used when
checking whether a proposed statement was genuinely absent or merely moved.

\section{Reproducibility notes}
\label{app:repro}

The archive accompanying this report contains:
\begin{itemize}[leftmargin=8mm]
\item \path{q_richardson_sinc_frontier.tex} --- the present source;
\item \path{q_richardson_sinc_frontier.pdf} --- the rendered report;
\item \path{verify_frontier_results.py} --- exact symbolic and high-precision
numerical checks;
\item \path{numerical_results.txt} and \path{numerical_results.csv} ---
machine-readable output;
\item \path{README.txt} --- compilation and reproduction commands.
\end{itemize}

The script checks the q-moment identity symbolically for several arbitrary
orders, generates the exact Bernoulli and Bell coefficients, verifies agreement
between the explicit and triangular filters, reproduces
\cref{tab:error1,tab:error3,tab:error10}, checks the leading constants, and
asserts the two-sided certificate \eqref{eq:certificate} at representative
points.

\begin{thebibliography}{99}

\bibitem{repo}
V. Reshetnikov,
\emph{ProveIt: Analysis/FabiusFunction documentation and Lean development},
GitHub repository,
\url{https://github.com/VladimirReshetnikov/ProveIt/tree/main/Analysis/FabiusFunction}.

\bibitem{Fabius1966}
J. Fabius,
\emph{A probabilistic example of a nowhere analytic \(C^\infty\)-function},
Z. Wahrscheinlichkeitstheorie verw. Gebiete 5 (1966), 173--174.

\bibitem{Arias1982}
J. Arias de Reyna,
\emph{An infinitely differentiable function with compact support: Definition
and properties},
Rev. Real Acad. Cienc. Exact. Fis. Natur. Madrid 76 (1982), 21--38;
English source at arXiv:1702.05442.

\bibitem{AriasArithmetic}
J. Arias de Reyna,
\emph{Arithmetic of the Fabius function},
arXiv:1702.06487 (2017).

\bibitem{Aistleitner}
C. Aistleitner, R. Hofer, and G. Larcher,
\emph{On parametric Thue--Morse sequences and lacunary trigonometric products},
arXiv:1502.06738; Ann. Inst. Fourier 67 (2017), 637--687.

\bibitem{Toth}
L. Toth,
\emph{Linear combinations of Dirichlet series associated with the Thue--Morse
sequence},
Integers 22 (2022), Article 98; arXiv:2211.13570.

\bibitem{Sidi}
A. Sidi,
\emph{Practical Extrapolation Methods: Theory and Applications},
Cambridge University Press, 2003.

\bibitem{Comtet}
L. Comtet,
\emph{Advanced Combinatorics: The Art of Finite and Infinite Expansions},
D. Reidel, 1974.

\bibitem{GasperRahman}
G. Gasper and M. Rahman,
\emph{Basic Hypergeometric Series}, 2nd ed.,
Cambridge University Press, 2004.

\end{thebibliography}

\end{document}
